%
% Copyright (c) 2016 Radoslaw Kujawa.
% All rights reserved.
%
% Redistribution and use in source and binary forms, with or without
% modification, are permitted provided that the following conditions
% are met:
%
% 1. Redistributions of source code must retain the above copyright
%    notice, this list of conditions and the following disclaimer.
% 2. Redistributions in binary form must reproduce the above copyright
%    notice, this list of conditions and the following disclaimer in the
%    documentation and/or other materials provided with the distribution.
%
% THIS SOFTWARE IS PROVIDED BY RADOSLAW KUJAWA (THE AUTHOR) AND CONTRIBUTORS
% ``AS IS'' AND ANY EXPRESS OR IMPLIED WARRANTIES, INCLUDING, BUT NOT LIMITED
% TO, THE IMPLIED WARRANTIES OF MERCHANTABILITY AND FITNESS FOR A PARTICULAR
% PURPOSE ARE DISCLAIMED.  IN NO EVENT SHALL THE AUTHOR OR CONTRIBUTORS
% BE LIABLE FOR ANY DIRECT, INDIRECT, INCIDENTAL, SPECIAL, EXEMPLARY, OR
% CONSEQUENTIAL DAMAGES (INCLUDING, BUT NOT LIMITED TO, PROCUREMENT OF
% SUBSTITUTE GOODS OR SERVICES; LOSS OF USE, DATA, OR PROFITS; OR BUSINESS
% INTERRUPTION) HOWEVER CAUSED AND ON ANY THEORY OF LIABILITY, WHETHER IN
% CONTRACT, STRICT LIABILITY, OR TORT (INCLUDING NEGLIGENCE OR OTHERWISE)
% ARISING IN ANY WAY OUT OF THE USE OF THIS SOFTWARE, EVEN IF ADVISED OF THE
% POSSIBILITY OF SUCH DAMAGE.
%
% 
\documentclass[dvipsnames,table]{beamer}
\usepackage{polski}

\usetheme{Rochester}
\usecolortheme{orchid}

\usepackage{listings}
\usepackage{ucs}
\usepackage[utf8x]{inputenc}
\usepackage{wasysym}
\usepackage[normalem]{ulem}
\usepackage{amsmath}
\usepackage{hyperref}
\usepackage{tikzsymbols}

\setbeamertemplate{navigation symbols}{}
\setbeamertemplate{caption}[numbered]
\setbeamerfont{caption}{size=\scriptsize}
\setbeamercolor{framenote}{bg=OSEC-red!25}
\setbeamercolor{rednote}{bg=Red!25}
\setbeamercolor{palette primary}{use=structure,fg=white,bg=OSEC-red}
\setbeamercolor{palette secondary}{use=structure,fg=white,bg=OSEC-red2}

\setbeamertemplate{itemize item}{\scriptsize\raise1pt\hbox{\donotcoloroutermaths$\blacktriangleright$}}
\setbeamertemplate{itemize subitem}{\tiny\raise1pt\hbox{\donotcoloroutermaths$\bullet$}}
\setbeamertemplate{itemize subsubitem}{\tiny\raise1pt\hbox{\donotcoloroutermaths{--}}}

\setbeamertemplate{enumerate item}{\insertenumlabel.}
\setbeamertemplate{enumerate subitem}{\insertenumlabel.\insertsubenumlabel}
\setbeamertemplate{enumerate subsubitem}{\insertenumlabel.\insertsubenumlabel.\insertsubsubenumlabel}
\setbeamertemplate{enumerate mini template}{\insertenumlabel}

\setbeamercolor{itemize item}{fg=OSEC-red, bg=OSEC-red}
\setbeamercolor{itemize subitem}{fg=OSEC-red, bg=OSEC-red}
\setbeamercolor{itemize subsubitem}{fg=OSEC-red, bg=OSEC-red}

\setbeamercolor{section number projected}{fg=white,bg=OSEC-red}
\setbeamercolor{subsection number projected}{fg=white,bg=OSEC-red}
\setbeamercolor{button}{bg=OSEC-red,fg=white}

\setbeamertemplate{section in toc}[circle]
\setbeamertemplate{subsection in toc}[square]

\definecolor{OSEC-red}{RGB}{160,29,44}
\definecolor{OSEC-red2}{RGB}{177,76,12}
\hypersetup{colorlinks=true,linkcolor=white,urlcolor=OSEC-red}

\setlength{\tabcolsep}{8pt}
\renewcommand{\arraystretch}{1.2}

\title{systemd-nspawn \\ kontenery bliższe sercu administratora}
\author{Radosław Kujawa -- radoslaw.kujawa@osec.pl}
\institute{OSEC}

\begin{document}

\begin{frame}
	\titlepage
\end{frame}

\begin{frame}[allowframebreaks]
	\frametitle{Spis treści}
	{
		\hypersetup{colorlinks=true,linkcolor=black,urlcolor=OSEC-red}
		\tableofcontents
	}
\end{frame}


\begin{frame}
\frametitle{Kontenery dziś}
\begin{itemize}
	\item Ogromne ułatwienie procesu developmentu oraz testowania.
	\item Potencjalne utrudnienie w procesie wdrożenia na produkcje oraz późniejszym zarządzaniu.
	\item Kontener - zunifikowany sposób dostawy aplikacji, od developera do administratora?
	\item Specyfikacja kontenera jako elementu budulcowego systemu dalej nie rozwiązuje problemów komunikacyjnych. 
\end{itemize}
\end{frame}

\begin{frame}
\frametitle{Kontenery systemd}
\begin{itemize}
	\item Oparte o mechanizm namespace.
	\item Dobrze zintegrowane z systemd.
	\item Chronione SELinuxem.
	\item Korzystające z narzędzi hosta do instalacji kontenerów, a także zarządzania nimi ({\tt machinectl}, {\tt systemctl} etc.).
	\item Domyślnie korzystają z systemu plików hosta, sieci hosta, oraz współdzielą namespace użytkowników (ale jest możliwość wydzielenia).
	\item Istnieje możliwość uruchamiania binarów innej architektury niż hosta (np. przez {\tt systemd-binfmt}).
\end{itemize}
\end{frame}

\begin{frame}
\frametitle{Instalacja}
\begin{itemize}
	\item Zalecane systemd 231 lub nowsze (Fedora 25).
	\item Paczka {\tt systemd-container} dostarcza potrzebne narzędzia.
	\item To wszystko! \Smiley
\end{itemize}
\end{frame}

\begin{frame}
\frametitle{Instalacja kontenera}
\begin{itemize}
	\item Fedora: {\tt dnf --releasever=25 --installroot=/container/foo/ install systemd passwd dnf fedora-release}
	\begin{itemize}
		\item Uwaga: cache {\tt dnf} jest {\it wewnątrz} kontenera... 
	\end{itemize}
	\item Debian: {\tt debootstrap --arch=amd64 stable /container/bar}
	\item Istnieje możliwość instalacji innej dystrybucji niż system hosta wewnątrz kontenera.
	\item Z wykorzystaniem innych narzędzi specyficznych dla danej dystrybucji.
\end{itemize}
\end{frame}

\begin{frame}
\frametitle{SELinux}
\begin{itemize}
	\item Typ {\tt container\_t} dla procesów oraz {\tt continer\_file\_t} i {\tt container\_ro\_file\_t}.
	\item {\tt semanage fcontext -a -t container\_file\_t '/container(/.*)?'}
	\item {\tt restorecon -R /container/}
	\item Możliwa bardziej zaawansowana polityka, wykraczająca poza typy {\tt container*}.
\end{itemize}	
\end{frame}

\begin{frame}
\frametitle{Start kontenera}
\begin{itemize}
	\item Ustawienie hasła roota
	\begin{itemize}
		\item {\tt systemd-nspawn -D /container/foo}
	\end{itemize}
	\item Start kontenera
	\begin{itemize}
		\item {\tt systemd-nspawn -bD /container/foo} (start interaktywny)
		\item {\tt systemd-nspawn -D /container/foo /opt/app} (uruchomienie polecenia w kontenerze)
		\item Tworzy jednostkę systemd {\tt machine-nazwa.scope} (transient).
		\item Jest także widoczny w {\tt machinectl -l}.
	\end{itemize}	
\end{itemize}
\end{frame}

\begin{frame}
\frametitle{Korzystanie z gotowych kontenerów}
\begin{itemize}
	\item Eksport własnego kontenera: {\tt machinectl export-tar kontener plik}.
	\item Import kontenera: {\tt machinectl ...}
	\item Import kontenera z sieci: {\tt machinectl pull-raw \href{https://dl.fedoraproject.org/pub/fedora/linux/releases/25/CloudImages/x86\_64/images/Fedora-Cloud-Base-25-1.3.x86\_64.raw.xz}{Obraz-Cloud-Fedora}}
	\item machinectl pull
\end{itemize}
\end{frame}


\begin{frame}
\frametitle{Zarządzanie z poziomu hosta}
\begin{itemize}
	\item {\tt machinectl status kontener}
	\item {\tt systemctl -M kontener}
	\item {\tt journalctl -M kontener}
	\item Procesy kontenera są widoczne w systemie hosta (dodatkowe pole {\tt machine} w {\tt ps}.
\end{itemize}
\end{frame}

\begin{frame}[fragile]
\frametitle{aaa}
\scriptsize
\begin{verbatim}
ooo oooooooo
oo
ooo
o
\end{verbatim}
\normalsize
\begin{itemize}
	\item foo {\tt bbb\_aaarr}.
\end{itemize}
\end{frame}

\begin{frame}
\frametitle{Zarządzanie storage kontenera}
\begin{itemize}
	\item System plików root - domyślnie osobny dla każdego kontenera
	\item Mount root ro
	\item btrfs
\end{itemize}
\end{frame}

\begin{frame}
\frametitle{Dostęp do katalogów hosta z kontenera}
\begin{itemize}
	\item --bind
\end{itemize}
\end{frame}

\begin{frame}
\frametitle{Zarządzanie siecią}
\begin{itemize}
	\item Wydzielenie interfejsu hosta dla kontenera: {\tt --network-interface}.
	\item Forward portów do kontenera {\tt --port}.
	\item Wirtualny interfejs między hostem a kontenerem: {\tt --network-veth}.
	\item Interfejs kontenera do bridge hosta: {\tt --network-bridge}
\end{itemize}
\end{frame}

\begin{frame}
\frametitle{Alternatywa dla Dockera?}
\begin{itemize}
	\item Mechanizmy dostarczane przez systemd są bardziej niskopoziomowe.
	\item {\tt rkt} używa {\tt systemd-nspawn}.
\end{itemize}	
\end{frame}

\begin{frame}
\frametitle{Podsumowując...}
\begin{itemize}
	\item Konteneryzacja oparta o {\tt systemd-nspawn} 
\end{itemize}
\end{frame}

\begin{frame}
\frametitle{Koniec\ldots}
\begin{center}
\includegraphics[scale=0.5]{img-oseclogo.png}

Dziękuje!

Czy są pytania?

\end{center}
\end{frame}
\end{document}

